% Generated from the matching Typst scaffold by
% scripts/migrate_typst_report_to_latex.py.

% ── 6.2  Interpretation and Implications ─────────────────────────────────
% PURPOSE: Turn the key findings into "when to forget", practitioner guidance, theory.
% MOVE: Pivot ("Rather than treat memory as always-on, treat forgetting as dosage").
% BEATS: see ordered bullets below — dosage → practitioner guidance → theory.
% FIGURES/TABLES: fig:interdependence (memory-lift by position + ~⅓ overlap).
% CITATIONS: xiong2025; anderson1991, richards2017 (transience as adaptive).
% SOURCE-DOC: README grounded-facts; docs/defense-qa.md (Q5).
% ─────────────────────────────────────────────────────────────────────────

\section{Interpretation and Implications}

\subsection{From always-on memory to footprint-free forgetting}

The dominant assumption in agent-memory engineering is that retention is an
unalloyed good: store every solved task, retrieve what is similar, and let the
accumulated experience compound. Our results invite a different default. With
retrieval held constant across all six policies (Invariant~\#5), no forgetting
policy is distinguishable from \texttt{full\_memory} on the CL-F1 resolve-rate
proxy --- every one of the five pre-registered Wilcoxon contrasts is
non-significant after Holm correction, and every rank-biserial effect carries a
bootstrap confidence interval that crosses zero. The benefit of unbounded
accumulation over carrying no memory at all is small (a $+1.6$~pp difference in
means, $0.445$ versus $0.429$; paired median $\Delta \approx +3.4$~pp) and itself
non-significant; the $N=8$ sequence-level design cannot resolve effects of this
magnitude, so we report it as a bounded benefit rather than a confirmed one.

What \emph{does} separate the policies is footprint. \texttt{full\_memory} carries
roughly $1{,}247$ tokens of injected context per task; the pruners carry between
$\sim\!600$ and $\sim\!700$ --- about half --- while \texttt{no\_memory} occupies
the zero-footprint corner (\Cref{fig:pareto-footprint}). Crucially, this saving is
\emph{not} a saving on compute: total tokens per run vary only from $5.74$M
(\texttt{no\_memory}) to $6.09$M (\texttt{full\_memory}), the five memory-carrying
policies fall within $\sim\!2\%$ of one another, and tool-call counts sit at
$\sim\!21$ for every policy. Cost in this agent is dominated by the
reasoning--action loop, not by the memory payload (\Cref{fig:pareto-compute}). The
practical implication reframes the question. Forgetting is not a lever that
recovers lost performance, nor a way to spend less compute; it is a way to keep the
injected context \emph{bounded at no measurable cost to the outcome}. Forgetting,
on this evidence, is close to free --- and ``free'' is precisely why it should be
the default rather than the exception.

\subsection{Dosage: target memory at the interdependent minority}

If accumulation is not what helps, the useful question is no longer ``store
everything or nothing'' but \emph{dosage}: how much memory, aimed at which tasks.
The mechanism behind the diluted average makes this concrete. Only a minority of
tasks structurally depend on a predecessor --- across the eight sequences the
fraction of tasks whose gold-patch files overlap an earlier task averages
$\sim\!0.3$, and ranges from $0.13$ (scikit-learn) to $0.57$ (pydata/xarray)
(\Cref{fig:interdependence}). This is a property of the benchmark, not of any
policy. For the file-disjoint majority, prior memory has nothing relevant to
contribute, so even a perfect retention policy is averaged against a large
population of tasks it cannot affect. ``Not every task needs memory'' is, on these
data, a measurable structural fact rather than a slogan.

Interdependence is necessary but not sufficient. Per-sequence memory lift (full
minus no-memory resolve rate) does \emph{not} track structural overlap: scikit-learn
has the lowest dependency ($0.13$) yet the highest lift ($+0.073$), while xarray has
the highest dependency ($0.57$) yet essentially zero lift, and matplotlib and
astropy show negative lift (\Cref{fig:memory-lift}). A dependency must additionally
be retrievable and the retrieved item must actually be used before it can flip an
outcome --- a chain of filters, each of which loses tasks. Forgetting is the
mechanism that realizes the dosage: by keeping memory small it concentrates the
limited retrieval budget on the few items most likely to matter, without paying the
footprint of carrying the rest. The contribution here is less a winning policy than
the quantification of \emph{why} the average benefit is small.

\subsection{Practitioner guidance}

Three recommendations follow directly. First, a \emph{simple bounded pruner is the
safe default}. A small memory budget (we used a cap of ten items) combined with a
content-agnostic eviction rule --- random or recency --- matches full accumulation
on the proxy at roughly half the footprint and at no compute penalty. There is no
reason to pay for unbounded storage growth that does not recover performance, and on
largely file-disjoint task streams an aggressive pruner is non-inferior; richer
retention should be reserved for repositories with demonstrably high task
interdependence.

Second, \emph{do not pay for content-awareness}. Neither content-aware policy beat
random pruning: CLS consolidation ($0.449$) is statistically indistinguishable from
random ($0.447$), and Type-Aware Decay is the lowest-scoring policy of all six
($0.427$). There is no evidence in these results that inspecting memory content to
decide what to keep buys anything over discarding items blindly.

Third, and most concretely, \emph{avoid pure retrieval-frequency decay} --- the
Type-Aware Decay trap. Its formal Wilcoxon contrast does not reach significance
(Holm $p = 1.000$), so this is a mechanistic caution rather than a measured loss,
but the warning signs converge: TAD carries the largest negative effect of any
policy (medium, $r_{rb} = -0.444$) and is the only policy whose mean falls below
\texttt{no\_memory}. The cause is structural. The retrieval-count term
$(1+\text{use\_count})^{0.5}$ in the decay formula (Invariant~\#8) overwhelms the
age term: early, high-use memories become effectively permanent while newly written
items are evicted almost immediately, so the store freezes at its first
$\sim\!10$ items and stops adapting. A decay rule that rewards past retrieval can
thus entrench stale experience --- precisely the experience-following risk
identified by \cite{xiong2025} --- and should be avoided in favour of simpler
bounded eviction.

\subsection{Where the field should look next}

The strongest forward-looking signal in our data is not about storage policy at all.
The task-level GLMM (Invariant~\#14) finds that none of the policy coefficients
reach significance, but two factors do: task difficulty and, decisively, sequence
position (coefficient $-1.49 \pm 0.086$). Agent effectiveness declines as a sequence
progresses regardless of which memory policy is in force, and accumulating memory
does not rescue it. This points away from eviction policy as the productive frontier
and toward the retrieval and trajectory pathway --- the position-dependent decline
that resembles context dilution rather than a failure of what to remember. If the
field wants to improve continual performance, the leverage appears to lie in how
relevant experience is surfaced and woven into a long trajectory, and in the
structure of task interdependence itself, not in the choice of forgetting rule. The
storage-policy question, at least under a held-constant retriever on this benchmark,
is largely settled in the negative; the retrieval-and-position question is open.

\subsection{Theoretical implication: transience as adaptive}

These findings sit comfortably with a long-standing view in cognitive science that
forgetting is not a defect of memory but a feature of it. The rational analysis of
memory holds that the decay of an item's availability tracks the declining
probability that it will be needed, so transience is an adaptation to the statistics
of the environment rather than a leak \cite{anderson1991}; more recent work argues
that active forgetting promotes generalisation and flexible behaviour rather than
degrading it \cite{richards2017}. Our footprint-non-inferiority result is the
engineering analogue: an agent that proactively discards most of its past pays no
measurable outcome penalty, because most of that past was never going to be
retrieved or used. Bounded, forgetful memory is not a compromise forced by resource
limits; on a task stream where only a minority of items remain relevant, it is the
rational policy.
